\section{Conversion Efficiency}
\label{sec:method-efficiency}

System efficiency was calculated from real AC input power and DC power at the
stated output measurement point:
\begin{equation}
 \eta_{\mathrm{sys}}=
 \frac{P_{\mathrm{out}}}{P_{\mathrm{in}}}\,100\%.
 \label{eq:efficiency-method}
\end{equation}
For steady voltage and current readings, output power is
\begin{equation}
 P_{\mathrm{out}}=V_{\mathrm{out}}I_{\mathrm{out}}.
 \label{eq:pout-method}
\end{equation}
For a changing waveform, the calculation instead needs the average of the
instantaneous product \(v(t)i(t)\).

For the V2 secondary-load tests, the AC supply's reported real power was
paired with voltage and current at the electronic-load terminals. Their
product gives the secondary-load power used in the efficiency estimate.
Loss in the leads between the secondary output and load is included with
the converter losses; no cable-loss correction was applied.

Paired entries use the same operating condition. DUT input measurements include
internal auxiliaries; the stated output sense point determines cable-loss
inclusion. No cable-loss correction was made without a measured voltage drop.
The results in Chapter~\ref{ch:results} are initial estimates because a
synchronized, repeated sweep with a complete uncertainty budget was not
available at every point. Unmatched observations were excluded.
